\chapter{Measure representation for BV functions}\label{sec:bv-representation-calculus}
This section proves the measure representation used in Definition \ref{def:bv} and proves constancy when the distributional derivative vanishes. The source definition and representation are \cite{Sim18}, Chapter 2, Section 2, equations (2.1)--(2.4), pp.~52--53. We retain the convention $\int u\,\operatorname{div}X=-\int X\cdot dDu$; changing the sign of a polar vector does not change its variation measure.

\begin{proposition}[Distributional derivatives as vector measures]\label{prop:bv-vector-representation}
Let $U\subset\mathbb R^n$ be open. For $u\in L^1_{\mathrm{loc}}(U)$, local finiteness of the variation supremum in Definition \ref{def:bv} is equivalent to representation of its distributional gradient by a locally finite vector-valued Radon measure. There is a Radon measure $\mu$ and a Borel vector field $\sigma$ such that
\[
 Du=\sigma\mu,\qquad |\sigma|=1\quad\mu\text{-almost everywhere},\qquad
 \mu(W)=\sup_{\substack{X\in C_c^1(W;\mathbb R^n)\\ |X|\leq1}}\int_Wu\,\operatorname{div}X
\]
for every open $W\subset U$. Thus $\mu=|Du|$ is exactly the total variation measure, including on open sets of infinite variation.
\end{proposition}
\begin{proof}
Define $L(X)=-\int u\,\operatorname{div}X$ on smooth compactly supported vector fields. The symmetry of the admissible tests under $X\mapsto-X$ bounds $|L(X)|$ by the local variation times $\|X\|_\infty$. Uniform approximation by smooth fields, with supports in a common compact neighborhood, extends $L$ uniquely to $C_c(U;\mathbb R^n)$. Independence of the approximation follows from the same local bound.

Apply Theorem \ref{thm:normed-vector-riesz} with the Euclidean norm on $\mathbb R^n$. It gives $L(X)=\int X\cdot\sigma\,d\mu$, with $|\sigma|=1$ almost everywhere. To identify the open-set supremum, note first that each admissible smooth test has $|L(X)|\leq\mu(W)$. Conversely, for $0\leq g\leq1$ compactly supported in $W$, the positive variation functional of Lemma \ref{lem:positive-vector-variation} gives
\[
 \int g\,d\mu=\sup\{L(F):F\in C_c(U;\mathbb R^n),\ |F|\leq g\}.
\]
Every such $F$ has support in $W$ and norm at most one. Extend it by zero and mollify with sufficiently small scale. Its smooth approximants stay supported in $W$, have norm at most one, and converge uniformly. Hence the original smooth supremum is at least $\int g\,d\mu$. Inner regularity and cutoff functions, exhausted over compact subsets of $W$, give the reverse inequality $\mu(W)$.

The coordinate distributional identities characterize $Du$ uniquely, since smooth tests determine Radon measures. Conversely, a locally finite vector measure representing $Du$ bounds the variation on $W\Subset U$ by its total variation there. The already established representation then identifies the minimal positive dominating measure with the stated supremum. This completes the equivalence without a separate vector-valued representation theorem left external.
\end{proof}

\begin{lemma}[A vanishing distributional gradient]\label{lem:zero-distributional-gradient}
If $U$ is connected and open and $u\in L^1_{\mathrm{loc}}(U)$ satisfies $Du=0$ in distributions, then $u$ is almost everywhere constant on $U$.
\end{lemma}
\begin{proof}
On a ball $B\Subset U$, convolve $u$ at sufficiently small scales. The distributional convolution identity gives zero classical gradient on any smaller concentric ball, so each mollification is constant there. It converges to $u$ in $L^1$ on that smaller ball. The constants form a Cauchy sequence, since the $L^1$ distance between two constants is the volume of the ball times their difference. Hence $u$ is almost everywhere constant on the ball. Constants on overlapping balls agree because the overlap has positive measure. A connected open subset of Euclidean space is joined by chains of overlapping balls; for example the set accessible from one ball by a finite chain is both relatively open and relatively closed. A countable ball cover removes any ambiguity from unions of null sets. The constants therefore agree almost everywhere on all of $U$.
\end{proof}
